\section{Experimental Evaluation}
\label{sec:experiments}

We evaluate \textsc{predictor4} on stratified samples from the ETP ground truth
\cite{etp2025} and compare it with \textsc{predictor3}.

\subsection{Setup}

\para{Ground truth.}
The ETP provides a complete $4{,}694 \times 4{,}694$ implication matrix for laws
of order at most 4, with all entries verified by Lean 4 proofs.  Positive
implications are encoded as $+3$ or $+4$ (two different verification methods);
non-implications as $-3$ or $-4$.  Of the $22{,}028{,}942$ off-diagonal entries,
$8{,}178{,}279$ ($37.1\%$) are implications and $13{,}855{,}357$ ($62.9\%$) are
non-implications.

\para{Test samples.}
We draw two stratified random samples.  The 50-case sample contains 25 implications
and 25 non-implications, balanced at the ground-truth ratio.  The 200-case sample
contains 100 implications and 100 non-implications.  Both samples are drawn with a
fixed random seed for reproducibility.

\para{Metrics.}
For each pair $(E, E')$, both classifiers return a confidence score $p \in (0, 1)$,
where $p > 0.5$ indicates implication and $p \leq 0.5$ indicates non-implication.
We report:
\begin{itemize}[leftmargin=*,itemsep=1pt,topsep=2pt]
  \item \textbf{Accuracy}: fraction of pairs correctly classified at threshold 0.5;
  \item \textbf{Average log-loss}: $\frac{1}{N}\sum_{i} \log p_i$ where $p_i$ is
        the confidence score assigned to the correct class (higher is better);
  \item \textbf{False positives (FP)}: non-implications classified as implications;
  \item \textbf{False negatives (FN)}: implications classified as non-implications.
\end{itemize}

\subsection{Main Results}

Table~\ref{tab:50case} reports results on the 50-case sample and
Table~\ref{tab:200case} on the 200-case sample.

\begin{table}[ht]
\centering
\caption{Results on the 50-case stratified sample (25 implications, 25 non-implications).
Best log-loss in \textbf{bold}.}
\label{tab:50case}
\renewcommand{\arraystretch}{1.1}
\begin{tabular}{@{}lcccc@{}}
\toprule
\textbf{Method} & \textbf{Accuracy} & \textbf{Avg.\ Log-Loss} & \textbf{FP} & \textbf{FN} \\
\midrule
\textsc{predictor3} & 50/50 (100\%) & \textbf{$-0.0001$} & 0 & 0 \\
\textsc{predictor4} & 50/50 (100\%) & $-0.0004$ & 0 & 0 \\
\bottomrule
\end{tabular}
\end{table}

\begin{table}[ht]
\centering
\caption{Results on the 200-case stratified sample (100 implications, 100 non-implications).
Runtime in seconds; per-case time in parentheses.}
\label{tab:200case}
\renewcommand{\arraystretch}{1.1}
\begin{tabular}{@{}lcccccc@{}}
\toprule
\textbf{Method} & \textbf{Accuracy} & \textbf{Avg.\ Log-Loss} & \textbf{FP} & \textbf{FN}
  & \textbf{Runtime (s)} \\
\midrule
\textsc{predictor3} & 200/200 (100\%) & $-0.0021$ & 0 & 0 & 93.4 (0.47/case) \\
\textsc{predictor4} & \textbf{200/200 (100\%)} & $-0.0025$ & 0 & 0 & 468.2 (2.34/case) \\
\midrule
\multicolumn{5}{l}{\textit{KB soundness:} 1/1 correct proofs (soundness 1.000)} \\
\bottomrule
\end{tabular}
\end{table}

Both classifiers achieve $100\%$ accuracy with zero false positives and zero false
negatives on both samples.  The slight increase in average log-loss for
\textsc{predictor4} ($-0.0003$/case on the 50-case sample, $-0.0004$/case on the
200-case sample) reflects cases where KB times out and the procedure falls through to
the BFS or structural estimation fallbacks, which return lower-confidence scores.

\subsection{KB Completion Analysis}

\para{Convergence rate.}
KB completion terminates (either successfully or with failure) within the 0.8-second
timeout for approximately $12\%$ of implication cases on the 200-case benchmark.
For the remaining $88\%$, KB times out and the BFS fallback is used.  The cases where
KB succeeds are those with simpler law structures (direct specializations, idempotency
consequences, and short derivations).

\para{Soundness verification.}
After enforcing the soundness guard, KB generates no false positives across all 200
tested cases (Theorem~\ref{thm:no-fp}).  KB successfully proved 1 implication on the
200-case sample; that proof was correct, yielding KB soundness~1.000.

\para{The invariant violation case.}
The specific case identified in Example~\ref{ex:bug} (law pair involving $x = y * ((x*y)*(z*z))$
and $x = (x*(x*y))*(z*w)$) is correctly handled after the fix: the buggy rule
$x \to y_0$ is rejected, KB returns inconclusive, and the negative tests correctly
classify the pair as a non-implication.

\subsection{Z3 Model Finding Analysis}

Z3 at domain sizes $n = 4$ and $n = 5$ is invoked only after the faster polynomial
checks and exhaustive size-2/3 enumeration have failed to find a counterexample.  On
the 200-case sample:
\begin{itemize}[leftmargin=*,itemsep=1pt,topsep=2pt]
  \item Most non-implications are caught by polynomial evaluation or exhaustive
        enumeration at sizes 2--3, typically within $0.1$~seconds.
  \item Z3 is needed for a minority of cases where smaller counterexamples do not exist.
  \item Every counterexample returned by Z3 was verified correct (no false positives
        from the Z3 component; Theorem~\ref{thm:z3-no-fp}).
\end{itemize}

\subsection{Runtime Analysis}

\textsc{predictor4} is approximately $5\times$ slower than \textsc{predictor3} on
the 200-case benchmark ($2.34$~s vs.\ $0.47$~s per case).  The overhead comes
primarily from the $0.8$-second KB timeout applied to cases where KB does not converge.
The per-case breakdown is approximately: $0.8$~s for KB timeout, $0.3$~s for Z3
at $n=4$, $0.6$~s for Z3 at $n=5$ (for cases reaching Z3), plus fast components.
This is a soundness-accuracy tradeoff: the KB-first approach introduces latency for
inconclusive cases but provides a sound proof for the $12\%$ where KB terminates.
